\section*{Chapter \ref{ch:b3:conformal} --- Conformal Maps and
the Riemann Mapping Theorem}

\begin{solution}{exo:b3:conformal:1}
(a) $\varphi$ and $\psi(w) = \iu\frac{1+w}{1-w}$ compose to
the identity in both orders (direct computation); $\varphi$
maps $\mathbb H$ into $\mathbb D$ and $\psi$ back
(\cref{ex:b3:conformal:mobius}). Values: $\varphi(\iu) = 0$,
$\varphi(0) = -1$, $\varphi(1) = \frac{1 - \iu}{1 + \iu} =
-\iu$, and $\varphi(z) \to 1$ as $z \to \infty$.

(b) Circles and lines are the solution sets of
$\alpha\abs z^2 + \bar\beta z + \beta\bar z + \gamma = 0$
($\alpha, \gamma \in \R$, $\beta \in \C$, $\abs\beta^2 >
\alpha\gamma$): $\alpha \neq 0$ circles, $\alpha = 0$ lines.
Substituting $z = 1/w$ and multiplying by $\abs w^2$:
$\gamma\abs w^2 + \beta w + \bar\beta\bar w + \alpha = 0$ ---
same family. Affine maps clearly preserve the family, and
every M\"obius map is a composition of affine maps and one
inversion ($\frac{az+b}{cz+d} = \frac ac + \frac{bc -
ad}{c}\cdot\frac1{cz + d}$ for $c \neq 0$).
\end{solution}

\begin{solution}{exo:b3:conformal:2}
(a) Schwarz gives $\abs{f(a)} \leq \abs a$ with equality (both
sides $= \abs a$): the equality case forces $f(z) =
\eu^{\iu\theta}z$, and $f(a) = a$ pins $\eu^{\iu\theta} = 1$.

(b) Let $a \neq b$ be fixed points and $g =
\varphi_a\circ f\circ\varphi_{-a} \in
\operatorname{Aut}(\mathbb D)$ --- using
\cref{thm:b3:conformal:autdisc} for $\varphi_{\pm a}$. Then
$g(0) = \varphi_a(f(a)) = 0$ and $g(c) = c$ for $c =
\varphi_a(b) \neq 0$: by (a), $g = \mathrm{id}$, so $f =
\varphi_{-a}\circ\varphi_a = \mathrm{id}$.
\end{solution}

\begin{solution}{exo:b3:conformal:3}
Fix $z$ and set $g = \varphi_{f(z)}\circ f\circ\varphi_{-z}$:
holomorphic $\mathbb D \to \mathbb D$ with $g(0) = 0$, so
$\abs{g'(0)} \leq 1$ (Schwarz). Chain rule with
$\varphi_a'(\zeta) = \frac{1 - \abs a^2}{(1 - \bar
a\zeta)^2}$:
\[
g'(0) = \varphi_{f(z)}'\bigl(f(z)\bigr)\cdot f'(z)\cdot
\varphi_{-z}'(0)
= \frac{1}{1 - \abs{f(z)}^2}\cdot f'(z)\cdot(1 - \abs z^2),
\]
whence the Schwarz--Pick inequality. Equality at some $z$
makes $g$ a rotation, hence $f =
\varphi_{-f(z)}\circ(\text{rotation})\circ\varphi_z \in
\operatorname{Aut}(\mathbb D)$ --- and then equality holds
everywhere (compute, or reapply with roles of $f, f^{-1}$
exchanged). Holomorphic self-maps of the disc are
$1$-Lipschitz for the hyperbolic metric $\frac{2\abs{\dd
z}}{1 - \abs z^2}$; automorphisms are its isometries.
\end{solution}

\begin{solution}{exo:b3:conformal:4}
(a) $z \mapsto \eu^z$: maps $\{0 < \operatorname{Im}z <
\pi\}$ bijectively onto $\mathbb H$ ($\eu^{x+\iu y} =
\eu^x\eu^{\iu y}$: modulus free, argument $y \in
\intoo0\pi$), holomorphic with nonvanishing derivative and
holomorphic inverse (principal $\log$).
(b) $z \mapsto z^2$ doubles arguments: the open quadrant
$\{0 < \arg z < \frac\pi2\}$ maps conformally onto $\mathbb
H$ (inverse: principal square root).
(c) $z \mapsto \frac{1 + z}{1 - z}$ maps $\mathbb D$ onto the
right half-plane and preserves the upper/lower symmetry: it
sends the upper half-disc onto the first quadrant; then
square, by (b), to reach $\mathbb H$: $z \mapsto \bigl(\frac{1
+ z}{1 - z}\bigr)^2$.
(d) The Blaschke factor $\varphi_{1/2}(z) = \frac{z -
\frac12}{1 - \frac z2}$: $\varphi_{1/2}(\tfrac12) = 0$ and
$\varphi_{1/2}'(\tfrac12) = \frac{1 - \frac14}{(1 -
\frac14)^2} = \frac43 > 0$.
\end{solution}

\begin{solution}{exo:b3:conformal:5}
(a) A conformal $\C \to \mathbb D$ (or $\mathbb H$, after
composing with Cayley) is a bounded entire function:
constant by Liouville --- not bijective. A conformal $\mathbb
D \to \C$ would have a conformal inverse $\C \to \mathbb D$:
same contradiction.

(b) The square is convex, hence simply connected, and proper:
conformally $\mathbb D$ (\cref{thm:b3:conformal:rmt}). The cut
plane $\C\setminus\intoc{-\infty}0$ is star-shaped about $1$
(segments from $1$ avoid the cut), hence simply connected,
and proper: conformally $\mathbb D$. The punctured disc: if
$g \colon \mathbb D\setminus\{0\} \to \mathbb D$ were
conformal, $g$ is bounded, so $0$ is removable
(\cref{thm:b3:residues:riemanncw}): $g$ extends to $\tilde g
\colon \mathbb D \to \mathbb D$, and $\tilde g(0)$, being in
the open image $g(\mathbb D\setminus\{0\}) = \mathbb D$, is
also $g(w)$ for some $w \neq 0$; two disjoint neighborhoods
of $0$ and $w$ have images that are open and share the value
$\tilde g(0)$, hence share \emph{other} values too (open
sets): $g$ takes some value twice on
$\mathbb D\setminus\{0\}$ --- contradicting injectivity. The
annulus $A = \{1 < \abs z < 2\}$: suppose $F \colon \mathbb D
\to A$ conformal. $F$ is zero-free on the simply connected
$\mathbb D$, so $F = \eu^L$ for holomorphic $L$
(\cref{def:b3:conformal:simplyconnected}). Let $\sigma$ be
the circle $\abs z = \frac32$ in $A$ and $\gamma =
F^{-1}\circ\sigma$, a closed path in $\mathbb D$; then
\[
1 = \operatorname{Ind}_\sigma(0)
= \frac1{2\iu\pi}\int_{F\circ\gamma}\frac{\dd w}{w}
= \frac1{2\iu\pi}\int_\gamma\frac{F'}{F}
= \frac1{2\iu\pi}\int_\gamma L' = 0
\]
($L'$ has a primitive): contradiction. Neither the punctured
disc nor the annulus is a disc in disguise.
\end{solution}

\begin{solution}{exo:b3:conformal:6}
(a) $T(w) = \frac{w - 1}{w + 1}$ maps $\{\operatorname{Re}w >
0\}$ conformally onto $\mathbb D$ (Cayley rotated: $\abs{w -
1} < \abs{w + 1}$ iff $\operatorname{Re}w > 0$), with $T(1) =
0$. For $f \in \mathcal F$, $g = T\circ f\colon \mathbb D \to
\mathbb D$ is holomorphic with $g(0) = 0$: Schwarz gives
$\abs{g(z)} \leq \abs z$.

(b) Inverting $T$: $f = \frac{1 + g}{1 - g}$, so
\[
\abs{f(z)} \leq \frac{1 + \abs{g(z)}}{1 - \abs{g(z)}}
\leq \frac{1 + \abs z}{1 - \abs z} :
\]
locally bounded (uniformly on $\abs z \leq r < 1$). Equality
at $z_0 \neq 0$ forces $\abs{g(z_0)} = \abs{z_0}$ and
alignment: $g$ a rotation, i.e.\ $f(z) = \frac{1 +
\eu^{\iu\theta}z}{1 - \eu^{\iu\theta}z}$ --- the
\emph{Herglotz extremals}, conformal maps onto the right
half-plane.
\end{solution}

\begin{solution}{exo:b3:conformal:7}
(i) Simple connectivity enters exactly twice, through the
existence of holomorphic square roots of zero-free functions
(\cref{def:b3:conformal:simplyconnected}): in Step 0 (the root
of $z - b$) and in Step 2 (the root of $\varphi_a\circ f$).
(ii) $\Omega \neq \C$ provides the point $b$ of Step 0 ---
without it the family $\mathcal F$ is empty of injective
bounded maps (Liouville). (iii) Connectedness is used whenever
the identity theorem or Hurwitz (\cref{exo:b3:residues:8})
speaks: the extremal limit is ``injective or constant'', and
constancy is excluded by $s > 0$; also in ``zero derivative
implies constant''. Failure for $\C$: Step 0 impossible, and
the conclusion is false (\cref{exo:b3:conformal:5}(a)).
Failure for the annulus: not simply connected --- the
square-root construction breaks (e.g.\ $z$ itself, zero-free
on $A$, has no holomorphic square root: the same index
computation as in \cref{exo:b3:conformal:5}(b) with
$\frac12\int_\sigma\frac{\dd z}z \notin 2\iu\pi\Z$) --- and
the conclusion is false too.
\end{solution}

\begin{solution}{exo:b3:conformal:8}
Substituting the Fourier expansion of $g$ into the Poisson
integral and using
$\frac1{2\pi}\int P_r(\varphi - t)\eu^{\iu nt}\dd t =
r^{\abs n}\eu^{\iu n\varphi}$ (read off $P_r$'s series):
$u(r\eu^{\iu\varphi}) = \sum_nc_n(g)\,r^{\abs
n}\eu^{\iu n\varphi}$, the interchange justified by normal
convergence ($\abs{c_n} \leq \norm g_\infty$, $r < 1$).

(a) $g = \cos t$: $c_{\pm1} = \frac12$, so $u =
r\cos\varphi = \operatorname{Re}z = x$ --- indeed harmonic
with the right boundary values.

(b) $\cos^2t = \frac12 + \frac{\cos 2t}2$: $u = \frac12 +
\frac{r^2\cos2\varphi}2 = \frac12 +
\frac{\operatorname{Re}(z^2)}2 = \frac12 + \frac{x^2 -
y^2}{2}$.

(c) $g = \mathbf 1_{(0,\pi)}$ (upper semicircle): $c_0 =
\frac12$ and $c_n = \frac{1 - (-1)^n}{2\iu\pi n}$ for $n \neq
0$, so
\[
u(r\eu^{\iu\varphi}) = \frac12 + \frac2\pi\sum_{k\geq0}
\frac{r^{2k+1}\sin\bigl((2k+1)\varphi\bigr)}{2k + 1},
\qquad u(0) = \frac12 :
\]
the center sees exactly the average of the boundary data ---
the mean value property in person.
\end{solution}

\begin{solution}{exo:b3:conformal:9}
From $(1-r)^2 \leq 1 - 2r\cos\theta + r^2 \leq (1+r)^2$:
\[
\frac{1-r}{1+r} = \frac{1 - r^2}{(1+r)^2} \leq P_r(\theta)
\leq \frac{1 - r^2}{(1 - r)^2} = \frac{1+r}{1-r} .
\]
For $u \geq 0$ harmonic on $\mathbb D$ and $s < 1$: $u_s(z) =
u(sz)$ is harmonic on a neighborhood of $\bar{\mathbb D}$,
hence equals its Poisson integral
(\cref{thm:b3:conformal:poisson}, uniqueness, applied to its
own boundary values); sandwiching the kernel and using the
mean value $\frac1{2\pi}\int u_s(\eu^{\iu t})\dd t = u(0)$:
\[
\frac{1-r}{1+r}\,u(0) \leq u(s\,r\eu^{\iu\varphi}) \leq
\frac{1+r}{1-r}\,u(0).
\]
Let $s \to 1^-$ at fixed $r\eu^{\iu\varphi}$ (continuity of
$u$): the Harnack inequalities. If $u$ is harmonic on $\C$
with $u \geq m$: apply Harnack to $u - m$ on discs $D(0, R)$,
i.e.\ to $z \mapsto u(Rz) - m$: for fixed $z$ and $r =
\abs z/R \to 0$, both bounds tend to $u(0) - m$: $u(z) =
u(0)$ --- constant (a two-sided Liouville from a one-sided
bound).
\end{solution}

\begin{solution}{exo:b3:conformal:10}
Locally, $u = \operatorname{Re}F$ with $F$ holomorphic
(\cref{prop:b3:conformal:harmonicholo}), so $u\circ f =
\operatorname{Re}(F\circ f)$ is harmonic wherever defined:
harmonicity is conformally invariant. On $\mathbb H$: the
principal logarithm gives $\log z = \ln\abs z + \iu\arg z$
holomorphic on $\mathbb H$, so $u = \frac1\pi\arg z =
\operatorname{Im}\bigl(\frac1\pi\log z\bigr)$ is harmonic,
with boundary limits: for $x > 0$, $\arg \to 0$, $u \to 0$;
for $x < 0$, $\arg \to \pi$, $u \to 1$: the data $\mathbf
1_{\intoo{-\infty}0}$ at every $x \neq 0$. Transporting by
the Cayley map (which sends $\mathbb D \to \mathbb H$ after
inversion and matches the upper semicircle to the negative
axis, up to the rotation fixed by chasing three boundary
points), $u\circ(\text{Cayley})$ solves the disc problem of
\cref{exo:b3:conformal:8}(c); evaluating at the center
retrieves $u = \frac12$ there, and the closed form
$\frac1\pi\arg$ can be checked against the series by summing
$\sum\frac{r^{2k+1}\sin((2k+1)\varphi)}{2k+1} =
\frac12\arctan\frac{2r\sin\varphi}{1 - r^2}$-type identities
--- the elementary route to the same answer.
\end{solution}

\begin{solution}{pb:b3:conformal:1}
\textbf{1.} $g$ is injective and holomorphic; on $C_\rho$
($\rho > 1$) it is smooth, and the enclosed area is
$\frac1{2\iu}\oint_{g(C_\rho)}\bar\zeta\,\dd\zeta$ (the Year 2
Green--Riemann area formula, applied with positive
orientation). Substituting $\zeta = g(w)$, $w =
\rho\eu^{\iu\theta}$:
\[
A_\rho = \frac1{2\iu}\int_{C_\rho}\overline{g(w)}\,g'(w)\dd
w = \frac\rho2\int_0^{2\pi}\overline{g(\rho\eu^{\iu\theta})}
\,g'(\rho\eu^{\iu\theta})\,\eu^{\iu\theta}\,\dd\theta .
\]
Insert $\bar g = \rho\eu^{-\iu\theta} + \bar b_0 +
\sum_m\bar b_m\rho^{-m}\eu^{\iu m\theta}$ and $g' = 1 -
\sum_nnb_n\rho^{-n-1}\eu^{-\iu(n+1)\theta}$: after
multiplying by $\eu^{\iu\theta}$, only frequency-zero
products survive the $\theta$-integration (normal
convergence justifies term-by-term work): the pair
$(\rho\eu^{-\iu\theta})\cdot1$ contributes $2\pi\rho$, and
each pair $\bar b_n\rho^{-n}\eu^{\iu
n\theta}\cdot(-nb_n\rho^{-n-1}\eu^{-\iu(n+1)\theta})$
contributes $-2\pi n\abs{b_n}^2\rho^{-2n-1}$:
\[
A_\rho = \pi\rho^2 - \pi\sum_{n\geq1}n\abs{b_n}^2\rho^{-2n}.
\]

\textbf{2.} Areas are nonnegative: $\sum_{n\leq
N}n\abs{b_n}^2\rho^{-2n} \leq \rho^2$ for every $N$; let
$\rho \to 1^+$ then $N \to \infty$:
$\sum_{n\geq1}n\abs{b_n}^2 \leq 1$. Equality in $\abs{b_1}
\leq 1$ forces all other $b_n = 0$: $g(w) = w + b_0 +
\eu^{\iu\alpha}/w$, which maps onto the complement of a
segment of length $4$ (a Joukowski-type map): the extremals.

\textbf{3.} $f(z)/z = 1 + a_2z + \cdots$ is holomorphic and
zero-free on $\mathbb D$ ($f$ vanishes only at $0$, simply:
injectivity), hence so is $f(z^2)/z^2$, which has a
holomorphic square root $\varphi$ with $\varphi(0) = 1$
(\cref{def:b3:conformal:simplyconnected}; $\mathbb D$ is
convex). Then $h(z) = z\varphi(z^2)$ satisfies $h^2 =
f(z^2)$, $h(z) = z\bigl(1 + \frac{a_2}2z^2 + \cdots\bigr)$
(binomial series for the root), and $h(z) = z\varphi(z^2)$ is
odd by construction ($\varphi(z^2)$ is even). Injectivity: $h(z) = h(z')$ gives $f(z^2) =
f(z'^2)$, so $z^2 = z'^2$, i.e.\ $z' = \pm z$; if $z' = -z$
then oddness gives $h(z) = -h(z)$, so $h(z) = 0$, forcing $z
= 0$ ($\varphi$ zero-free): $z = z' = 0$.

\textbf{4.} $g(w) = 1/h(1/w)$: for $\abs w > 1$, $1/w \in
\mathbb D\setminus\{0\}$ and $h \neq 0$ there: well defined,
injective (composition of injections), with expansion
\[
g(w) = \frac{1}{\frac1w\bigl(1 + \frac{a_2}{2w^2} +
\cdots\bigr)} = w\Bigl(1 - \frac{a_2}{2w^2} + \cdots\Bigr)
= w - \frac{a_2}{2}\,w^{-1} + \cdots :
\]
of the Part I form with $b_1 = -\frac{a_2}2$. The area theorem
gives $\abs{\frac{a_2}2} \leq 1$: $\abs{a_2} \leq 2$.

\textbf{5.} $k(z) = \frac z{(1-z)^2} = z\sum_{m\geq0}(m +
1)z^m = \sum_{n\geq1}nz^n$: coefficients $a_n = n$, so $a_2 =
2$. Univalence and image: $k = \frac14\bigl[w^2 - 1\bigr]$
with $w = \frac{1 + z}{1 - z}$, a conformal map of $\mathbb D$
onto the right half-plane; $w^2$ maps that half-plane
conformally onto $\C\setminus\intoc{-\infty}0$; then
$\frac{\cdot - 1}4$ gives
$\C\setminus\intoc{-\infty}{-\frac14}$: injective at each
stage, image as claimed.

\textbf{6.} $F = \frac{cf}{c - f}$: since $c \notin f(\mathbb
D)$, the denominator never vanishes: $F$ is holomorphic, and
injective ($w \mapsto \frac{cw}{c - w}$ is M\"obius, injective
off $w = c$). Expansion: with $f = z + a_2z^2 + \cdots$,
\[
F = f\cdot\frac1{1 - f/c} = \bigl(z + a_2z^2\bigr)\Bigl(1 +
\frac zc\Bigr) + O(z^3)
= z + \Bigl(a_2 + \frac1c\Bigr)z^2 + O(z^3):
\]
$F \in \mathcal S$ with $A_2 = a_2 + \frac1c$.

\textbf{7.} Bieberbach twice: $\abs{a_2} \leq 2$ and
$\abs{a_2 + \frac1c} \leq 2$, so $\abs{\frac1c} \leq 4$:
$\abs c \geq \frac14$. Every omitted value lies outside
$D(0,\frac14)$, i.e.\ $f(\mathbb D) \supseteq D(0, \frac14)$.
Sharp: the Koebe function omits $-\frac14$ (question 5).

\textbf{8.} Let $d = d(z_0, \partial\Omega)$ and $\psi =
\varphi^{-1} \colon \mathbb D \to \Omega$, $\psi(0) = z_0$,
$\psi'(0) = 1/\varphi'(z_0) > 0$. The normalization $\tilde
f(z) = \frac{\psi(z) - z_0}{\psi'(0)}$ lies in $\mathcal S$,
so its image contains $D(0, \frac14)$; scaling back,
$\Omega = \psi(\mathbb D) \supseteq D\bigl(z_0,
\tfrac{\psi'(0)}4\bigr)$, so $d \geq \frac{\psi'(0)}4 =
\frac1{4\varphi'(z_0)}$: the left inequality. For the right: $\varphi\circ(z_0 + d\,\cdot)$
maps $\mathbb D$ into $\mathbb D$ with $0 \mapsto 0$; Schwarz
bounds its derivative at $0$: $d\,\varphi'(z_0) \leq 1$, i.e.\
$\frac1{\varphi'(z_0)} \geq d$ --- together
\[
d \leq \frac1{\varphi'(z_0)} \leq 4d .
\]
(The left inequality as displayed in the statement is the
same chain rearranged.)

\textbf{9.} For the Koebe function $a_n = n$: equality
throughout the conjecture; its rotations
$\eu^{-\iu\theta}k(\eu^{\iu\theta}z) = \sum
n\eu^{\iu(n-1)\theta}z^n$ have $\abs{a_n} = n$ too. Pushing
question 4: $h = z + \frac{a_2}2z^3 + c_5z^5 + \cdots$ with
$h^2 = f(z^2)$; comparing $z^6$-coefficients: $2c_5 +
\frac{a_2^2}4 = a_3$, so $c_5 = \frac{a_3}2 -
\frac{a_2^2}8$. Then
\[
g(w) = \frac1{h(1/w)} = w\Bigl(1 - \frac{a_2}{2w^2} +
\Bigl(\frac{a_2^2}4 - c_5\Bigr)\frac1{w^4} + \cdots\Bigr)
= w - \frac{a_2}2w^{-1} + \Bigl(\frac{3a_2^2}8 -
\frac{a_3}2\Bigr)w^{-3} + \cdots,
\]
and the area theorem ($\sum n\abs{b_n}^2 \leq 1$) yields
\[
\Bigl|\frac{a_2}2\Bigr|^2 + 3\,\Bigl|\frac{3a_2^2}8 -
\frac{a_3}2\Bigr|^2 \leq 1 .
\]
Koebe check ($a_2 = 2$, $a_3 = 3$): $\abs{\frac{a_2}2}^2 = 1$
and $\frac{3\cdot4}8 - \frac32 = 0$: total exactly $1$ ---
extremal, as it must be.
\end{solution}
